\chapter*{Executive Summary}
\addcontentsline{toc}{chapter}{Executive Summary}
\markboth{Executive Summary}{Executive Summary}

\textbf{BuildScan~3D} is a mobile-to-cloud system that allows users to turn a short video of an indoor space captured with a smartphone into an accurate 3D model that can be used to measure.

\paragraph{Problem.} The 3D scanning of interiors is usually performed using costly methods that involve either using LiDAR devices which cost more than \$10,000 and require multiple preconditions or using photogrammetry, which requires advanced experience in the sphere as well as much time to analyze the data. Many construction companies, facility managers, and students cannot afford to buy the necessary equipment and utilize modern technology to create highly accurate models.

\paragraph{Research.} We analyzed the existing developments related to the application of Gaussian Splatting for 3D scanning~\cite{kerbl3dgs}, PocketGS~\cite{pocketgs}, and visual-inertial SLAM (VIGS-SLAM~\cite{vigsslam}) and found a significant gap. Thus, no system has combined the use of smartphones equipped with ARCore and cloud computing, allowing receiving accurate 3D models in a fully automatic manner.

\paragraph{Solution.} Our pipeline includes five stages, which include: (1) capturing video at a rate of 30 fps, ARCore poses, depth and IMU data; (2) alignment of the captures; (3) processing with the use of COLMAP and Nerfstudio; (4) a structural layer (wall planes, corners, wall lengths, door detection) for floor plans and 1:1 measurement, prototype built; (5) viewing the data in a planned web viewer with metric measurements.

\paragraph{Current status (Iteration~1: core pipeline complete, refinement ongoing).} This Android capture app is capable of recording all necessary sensor data together with live IQA image quality processing and coverage guidance (using red/yellow/green dot overlay). FastAPI processing server receives the data and dispatches the jobs to RunPod GPU pods that are available as per request (the default pod is NVIDIA RTX 4090, but it can be selected at run-time). The approximate time taken in end-to-end processing using the development RTX 3070 has been calculated at about 75\,min; however, the cloud-GPU timings are yet to be measured, and the under-35-minute RTX 4090 figure is an FYP-2 projection. The current best reconstruction results in 26.52\,dB PSNR / 0.8886 SSIM for 42 frames that have not been used for training any of the models so far. Tape-measure accuracy is showing the mean error to be at 4.8\,mm (0.19\%) with the maximum value of 0.29\% (8\,mm in the case of 2.8\,m wall), which is well inside the target of 5\% accuracy. Python code, Android builds, and web UI are all verified via CI/CD process via GitHub Actions. Open items include refinement of 3DGS, creation of 3D mesh, completion of the structural layer, and web workspace.

\paragraph{Risks.} The cloud GPU availability and per-job cost (not yet measured); structural layer durability among various rooms (the ceiling-height detection is still not reliable); the success of measurement verification for various room layouts; the unmeasured RTX 4090 processing time. The complete risk register comprises 12 items with probability, effect, and mitigation involved.

\paragraph{FYP-2 plan.} Finish up the structural layer and automatic floor plans (Stage 4) and the web-based measurement environment (Stage 5); conduct field trials in further rooms and check correctness of measurement against ground data acquired by means of tape measure. 3D mesh building is also included in the plans yet not scheduled.
